\section{The coprime catalogue: $\gcd(E,o)=1$, $E+o\le 20$}

When $\gcd(E,o)=1$ we have $d=1$, so $R_1=\Z[t,t^{-1}]$ and \emph{every cyclic word is a
single Laurent polynomial in one variable}. Concretely the closure $\chi(w)=1$ is imposed by
the substitution of Proposition~\ref{prop:reparam} with $s=1$,
\[
   X\longmapsto t^{-(E-o)},\qquad Y\longmapsto t^{\,o},
\]
under which a monomial $X^{\alpha}Y^{\beta}$ of $S(w)$ becomes $t^{\,\beta o-\alpha(E-o)}$.
Clearing the lowest power of $t$ gives an integer polynomial $P_w(t)\in\Z[t]$, well defined up
to $t\mapsto t^{-1}$ and a power of $t$ (Proposition~\ref{prop:conj}); e.g.\ $uuv\mapsto t+1$
and $uuuv\mapsto t^2+t+1=\Phi_3$. The relevant invariants are those of $P_w$: its factorisation
into cyclotomic and non-cyclotomic parts, its Mahler measure, and the Galois groups of its
irreducible factors. (For general $d$ one keeps $s$: e.g.\ $uvuv\mapsto 1+s$ in $R_2$ and
$uuu\mapsto1+s+s^2$ in $R_3$; all $10$ classes with $E+o\le6$, any $d$, are cyclotomic.)

We enumerate \emph{every} coprime cyclic class with $E+o\le20$: $1233$ classes across $64$
types. All data are produced by \texttt{gen\_table.py} (the cocycle law of
Proposition~\ref{prop:cocycle} is re-verified on every split of every word up to length~$6$)
and recorded in full in \texttt{coprime\_catalogue.csv} (which lists each $P_w$ and its
Littlewood core $R_w$ of Proposition~\ref{prop:phio} explicitly); Table~\ref{tab:summary}
aggregates by type and Table~\ref{tab:highlights} lists the arithmetically extreme classes.

\subsection*{The four known cycles}

The four known cycles of the $3n+1$ map on $\Z$ are short words of the accelerated $T$-map
($u:x\mapsto(3x+1)/2$, $v:x\mapsto x/2$): the parity sequence of a periodic orbit is a positive
word $w$, its fixed point $F_w(x)=x$ is the cycle value, and the type $(E,o)$ gives
$D=2^{E}-3^{o}$. They are $\mathtt{u}$ ($\{-1\}$), $\mathtt{uv}$ ($\{1,2\}$), $\mathtt{uuv}$
(the $-5$ cycle), and $\mathtt{uuuuvuuuvvv}$ (the $-17$ cycle, $7$ odd steps); all four are
coprime. The Mahler measures and $D=2^E-3^o$ are exact; the fixed points reproduce the cycle
values (checked in \texttt{gen\_table.py}).

\begin{prop}[The cyclic factor and the Littlewood core]\label{prop:phio}
Let $w$ be a cyclic word of coprime type. Then $P_w$ is a $0/1$ polynomial with exactly $o$
terms, the full cyclotomic factor
\[
   g_o(t)\;:=\;\frac{t^{o}-1}{t-1}\;=\;1+t+\cdots+t^{o-1}\;=\!\!\prod_{d\mid o,\ d>1}\!\!\Phi_d(t)
\]
divides $P_w(t)$, and the quotient $R_w:=P_w/g_o$ has all coefficients in $\{-1,0,1\}$.
\end{prop}

\begin{proof}
By Proposition~\ref{prop:closed} (the $k$-th $u$ has $\alpha_k=k$ earlier $u$'s) the exponents
$e_k=o\beta_k-(E-o)k$ satisfy $e_k\equiv-(E-o)k\pmod o$; since $\gcd(E-o,o)=\gcd(E,o)=1$ the map
$k\mapsto-(E-o)k$ is a bijection of $\Z/o\Z$, so $\{e_k\bmod o\}$ is a complete residue system.
In particular the $e_k$ are pairwise distinct, so $P_w=\sum_k t^{e_k}$ is $0/1$ with $o$ terms;
and reducing each $t^{e_k}$ modulo $t^{o}-1$ gives $P_w\equiv\sum_{j=0}^{o-1}t^{j}=g_o\pmod{t^{o}-1}$,
so $P_w$ vanishes at every $o$-th root of unity except $1$, i.e.\ $g_o\mid P_w$.

For the quotient, write $[P]_m$ for the coefficient of $t^m$ and let $e(\rho)$ be the unique
exponent of $P_w$ in the class $\rho\bmod o$. From $R_w\,(t^{o}-1)=P_w\,(t-1)$ one gets
$R_n=R_{n+o}+[P_w(t-1)]_{n+o}$; summing downward from $R_{+\infty}=0$,
\[
   R_n=\sum_{i\ge1}\bigl([P_w]_{n+oi-1}-[P_w]_{n+oi}\bigr)
      =\mathbf 1\!\left[e\bigl((n{-}1)\bmod o\bigr)>n{-}1\right]
      -\mathbf 1\!\left[e\bigl(n\bmod o\bigr)>n\right],
\]
a difference of two $0/1$ indicators (one exponent per residue class), hence in $\{-1,0,1\}$.
\end{proof}

\begin{cor}[Normalisation]\label{cor:Rone}
$R_w(1)=1$; equivalently $R_w=1+(t-1)\,Q_w$, where
$Q_w(t)=\sum_{n\ge0}\mathbf 1\!\left[e(n\bmod o)>n\right]t^{n}$ is itself a $0/1$ polynomial. Hence
\[
   P_w \;=\; g_o\,\bigl(1+(t-1)\,Q_w\bigr),\qquad g_o,\,Q_w\in\{0,1\}[t],\quad R_w\in\{-1,0,1\}[t].
\]
\end{cor}

\begin{proof}
$P_w$ has exactly $o$ terms, so $R_w(1)=P_w(1)/g_o(1)=o/o=1$. With
$f(n):=\mathbf 1[e(n\bmod o)>n]$ the proof of Proposition~\ref{prop:phio} gives $R_n=f(n-1)-f(n)$,
whose tail sums telescope: $\bigl[(R_w-1)/(t-1)\bigr]_n=\sum_{i>n}R_i=f(n)\in\{0,1\}$.
\end{proof}

\begin{table}[h]\centering\small\renewcommand{\arraystretch}{1.25}
\begin{tabular}{@{}llccc@{}}
\toprule
word $w$ & cycle & $(E,o)$ & $D=2^{E}-3^{o}$ & $\mathrm{M}(P_w)$\\
\midrule
$\mathtt{u}$ & $\{-1\}$ & $(1,1)$ & $-1$ & $1$\\
$\mathtt{uv}$ & $\{1\}$ & $(2,1)$ & $1$ & $1$\\
$\mathtt{uuv}$ & $\{-5\}$ & $(3,2)$ & $-1$ & $1$\\
$\mathtt{uuuuvuuuvvv}$ & $\{-17\}$ & $(11,7)$ & $-139$ & $1.42747$\\
\bottomrule
\end{tabular}
\caption{The four known cycles of the $3n+1$ map on $\mathbb{Z}$, read as words of the accelerated $T$-map ($u:x\mapsto(3x+1)/2$, $v:x\mapsto x/2$): the cycle value $x$ (fixed point of the word map), the type, $D=2^E-3^o$, and the Mahler measure of $P_w$.}\label{tab:cycles}
\end{table}
\begin{align*}
P_{-5}(t) &= t + 1\\
 &= \Phi_{2}.\\[4pt]
P_{-17}(t) &= t^{17} + t^{13} + t^{9} + t^{8} + t^{5} + t^{4} + 1\\
 &= \Phi_{7}\,\bigl(t^{11} - t^{10} + t^{7} - t^{6} + t^{4} - t + 1\bigr).
\end{align*}


The two trivial cycles give $P_w=1$; the $-5$ cycle gives $P_w=\Phi_2$ (purely cyclotomic);
and the $-17$ cycle is the first to carry arithmetic, $P_w=g_7\,R_w$ with $g_7=\Phi_7$ (here
$o=7$ is prime) and Littlewood core $R_w=t^{11}-t^{10}+t^{7}-t^{6}+t^{4}-t+1$, irreducible and
non-reciprocal, of Mahler measure $1.42747\ldots$ (six roots off the unit circle). The
cyclotomic factor $g_o$ is forced by Proposition~\ref{prop:phio}; the core $R_w=P_w/g_o$ is the
genuine invariant attached to the cycle (Corollary~\ref{cor:Rone} writes it as $1+(t-1)Q_w$ with
$Q_w=t^{10}+t^{6}+t^{3}+t^{2}+t$).

\begin{prop}[Coprime cyclotomic frontier]\label{prop:cofrontier}
Every coprime cyclic word with $E+o\le 9$ has $\Mahler(P_w)=1$; i.e.\ $P_w$ is a product of
cyclotomic polynomials. The first non-cyclotomic class occurs at $E+o=10$, type $(7,3)$: for
$w=\mathtt{uuvuvvv}$,
\[
   P_w \;=\; \Phi_3(t)\,\bigl(t^{3}-t^{2}+1\bigr),\qquad
   \Mahler(P_w)=1.324718\ldots,
\]
the plastic number (the smallest Pisot number, real root of $x^{3}-x-1$); its non-cyclotomic
factor has Galois group $S_3$.
\end{prop}

\begin{prop}[Lehmer's number is realised]\label{prop:lehmer}
Across the $1233$ classes the Mahler measures lie in $[1,\,1.65378]$, with $274$ cyclotomic and
$959$ non-cyclotomic. The smallest measure exceeding $1$ is \emph{Lehmer's number}
$\Mahler=1.1762808\ldots$, attained (among others) by $w=\mathtt{uuuvuuvvvuuv}$ of type
$(12,7)$, whose non-cyclotomic factor is Lehmer's polynomial
\[
   t^{10}-t^{9}+t^{7}-t^{6}+t^{5}-t^{4}+t^{3}-t+1 .
\]
This is the smallest Mahler measure $>1$ known for any integer polynomial (Lehmer's problem);
that it is realised by a short Syracuse word is the catalogue's most striking feature.
\end{prop}

\begin{table}[h]
\centering\footnotesize
\renewcommand{\arraystretch}{1.1}
\setlength{\tabcolsep}{4.5pt}
\begin{tabular}{@{}cccccc@{\qquad}cccccc@{}}
\toprule
$(E,o)$ & cl. & cy. & nc. & min & max & $(E,o)$ & cl. & cy. & nc. & min & max\\
\midrule
$(1,1)$ & 1 & 1 & 0 & $--$ & $1.0000$ & $(11,1)$ & 1 & 1 & 0 & $--$ & $1.0000$ \\
$(2,1)$ & 1 & 1 & 0 & $--$ & $1.0000$ & $(11,2)$ & 5 & 5 & 0 & $--$ & $1.0000$ \\
$(3,1)$ & 1 & 1 & 0 & $--$ & $1.0000$ & $(11,3)$ & 15 & 5 & 10 & $1.3247$ & $1.3716$ \\
$(3,2)$ & 1 & 1 & 0 & $--$ & $1.0000$ & $(11,4)$ & 30 & 10 & 20 & $1.3247$ & $1.4864$ \\
$(4,1)$ & 1 & 1 & 0 & $--$ & $1.0000$ & $(11,5)$ & 42 & 8 & 34 & $1.2612$ & $1.5367$ \\
$(4,3)$ & 1 & 1 & 0 & $--$ & $1.0000$ & $(11,6)$ & 42 & 8 & 34 & $1.2612$ & $1.5367$ \\
$(5,1)$ & 1 & 1 & 0 & $--$ & $1.0000$ & $(11,7)$ & 30 & 10 & 20 & $1.3247$ & $1.4864$ \\
$(5,2)$ & 2 & 2 & 0 & $--$ & $1.0000$ & $(11,8)$ & 15 & 5 & 10 & $1.3247$ & $1.3716$ \\
$(5,3)$ & 2 & 2 & 0 & $--$ & $1.0000$ & $(11,9)$ & 5 & 5 & 0 & $--$ & $1.0000$ \\
$(5,4)$ & 1 & 1 & 0 & $--$ & $1.0000$ & $(12,1)$ & 1 & 1 & 0 & $--$ & $1.0000$ \\
$(6,1)$ & 1 & 1 & 0 & $--$ & $1.0000$ & $(12,5)$ & 66 & 5 & 61 & $1.1763$ & $1.5586$ \\
$(6,5)$ & 1 & 1 & 0 & $--$ & $1.0000$ & $(12,7)$ & 66 & 5 & 61 & $1.1763$ & $1.5586$ \\
$(7,1)$ & 1 & 1 & 0 & $--$ & $1.0000$ & $(13,1)$ & 1 & 1 & 0 & $--$ & $1.0000$ \\
$(7,2)$ & 3 & 3 & 0 & $--$ & $1.0000$ & $(13,2)$ & 6 & 6 & 0 & $--$ & $1.0000$ \\
$(7,3)$ & 5 & 3 & 2 & $1.3247$ & $1.3247$ & $(13,3)$ & 22 & 6 & 16 & $1.3247$ & $1.3756$ \\
$(7,4)$ & 5 & 3 & 2 & $1.3247$ & $1.3247$ & $(13,4)$ & 55 & 15 & 40 & $1.3247$ & $1.5189$ \\
$(7,5)$ & 3 & 3 & 0 & $--$ & $1.0000$ & $(13,5)$ & 99 & 10 & 89 & $1.1763$ & $1.6001$ \\
$(7,6)$ & 1 & 1 & 0 & $--$ & $1.0000$ & $(13,6)$ & 132 & 14 & 118 & $1.2353$ & $1.6538$ \\
$(8,1)$ & 1 & 1 & 0 & $--$ & $1.0000$ & $(13,7)$ & 132 & 14 & 118 & $1.2353$ & $1.6538$ \\
$(8,3)$ & 7 & 3 & 4 & $1.3247$ & $1.3497$ & $(14,1)$ & 1 & 1 & 0 & $--$ & $1.0000$ \\
$(8,5)$ & 7 & 3 & 4 & $1.3247$ & $1.3497$ & $(14,3)$ & 26 & 6 & 20 & $1.3247$ & $1.3768$ \\
$(8,7)$ & 1 & 1 & 0 & $--$ & $1.0000$ & $(14,5)$ & 143 & 9 & 134 & $1.2164$ & $1.6038$ \\
$(9,1)$ & 1 & 1 & 0 & $--$ & $1.0000$ & $(15,1)$ & 1 & 1 & 0 & $--$ & $1.0000$ \\
$(9,2)$ & 4 & 4 & 0 & $--$ & $1.0000$ & $(15,2)$ & 7 & 7 & 0 & $--$ & $1.0000$ \\
$(9,4)$ & 14 & 6 & 8 & $1.3247$ & $1.4458$ & $(15,4)$ & 91 & 21 & 70 & $1.3247$ & $1.5298$ \\
$(9,5)$ & 14 & 6 & 8 & $1.3247$ & $1.4458$ & $(16,1)$ & 1 & 1 & 0 & $--$ & $1.0000$ \\
$(9,7)$ & 4 & 4 & 0 & $--$ & $1.0000$ & $(16,3)$ & 35 & 7 & 28 & $1.3247$ & $1.3782$ \\
$(9,8)$ & 1 & 1 & 0 & $--$ & $1.0000$ & $(17,1)$ & 1 & 1 & 0 & $--$ & $1.0000$ \\
$(10,1)$ & 1 & 1 & 0 & $--$ & $1.0000$ & $(17,2)$ & 8 & 8 & 0 & $--$ & $1.0000$ \\
$(10,3)$ & 12 & 4 & 8 & $1.3247$ & $1.3679$ & $(17,3)$ & 40 & 8 & 32 & $1.3247$ & $1.3787$ \\
$(10,7)$ & 12 & 4 & 8 & $1.3247$ & $1.3679$ & $(18,1)$ & 1 & 1 & 0 & $--$ & $1.0000$ \\
$(10,9)$ & 1 & 1 & 0 & $--$ & $1.0000$ & $(19,1)$ & 1 & 1 & 0 & $--$ & $1.0000$ \\
\bottomrule
\end{tabular}
\caption{All $64$ coprime types ($\gcd(E,o)=1$, $E+o\le20$), in two side-by-side blocks
of $32$: number of cyclic classes (cl.), how many are cyclotomic (cy.) and non-cyclotomic
(nc.), the smallest Mahler measure $>1$ in the type (min; ``--'' if none), and the largest
Mahler measure (max).}\label{tab:summary}
\end{table}


\begin{table}[h]
\centering\small
\renewcommand{\arraystretch}{1.2}
\begin{tabular}{@{}llcccl@{}}
\toprule
$w$ & $(E,o)$ & $\deg P_w$ & $\deg_{\mathrm{nc}}$ & $\mathrm{M}$ & class\\
\midrule
$\mathtt{uuuvuuvvvuuv}$ & $(12,7)$ & $16$ & $10$ & $1.17628$ & Lehmer \\
$\mathtt{uuuvvuvvvuvv}$ & $(12,5)$ & $14$ & $10$ & $1.17628$ & Lehmer \\
$\mathtt{uuvuvvvuvvvuv}$ & $(13,5)$ & $14$ & $10$ & $1.17628$ & Lehmer \\
$\mathtt{uuuvvuvvvvvuvv}$ & $(14,5)$ & $18$ & $10$ & $1.21639$ & Salem \\
$\mathtt{uuuuvvuvvvuvv}$ & $(13,6)$ & $21$ & $16$ & $1.23526$ & 2 ext. \\
$\mathtt{uuuvuuvvvvuuv}$ & $(13,7)$ & $22$ & $16$ & $1.23526$ & 2 ext. \\
$\mathtt{uuvuuvvvuvvv}$ & $(12,5)$ & $16$ & $12$ & $1.25105$ & 2 ext. \\
$\mathtt{uuuvvvuvvuvvv}$ & $(13,5)$ & $16$ & $12$ & $1.25105$ & 2 ext. \\
$\mathtt{uuvvuvvuuvvvvv}$ & $(14,5)$ & $16$ & $12$ & $1.25105$ & 2 ext. \\
$\mathtt{uuuvuuuvvuvv}$ & $(12,7)$ & $18$ & $12$ & $1.25105$ & 2 ext. \\
$\mathtt{uuuuvuvvvuv}$ & $(11,6)$ & $15$ & $10$ & $1.26123$ & Salem \\
$\mathtt{uuuvuvvvvuv}$ & $(11,5)$ & $14$ & $10$ & $1.26123$ & Salem \\
$\mathtt{uuvuuvvvvuvvvv}$ & $(14,5)$ & $22$ & $18$ & $1.26633$ & 3 ext. \\
$\mathtt{uuuvuvuuuvvv}$ & $(12,7)$ & $16$ & $10$ & $1.28358$ & 2 ext. \\
$\mathtt{uuvuuvuuvvuvv}$ & $(13,7)$ & $16$ & $10$ & $1.28358$ & 2 ext. \\
$\mathtt{uuuuvvuvuvv}$ & $(11,6)$ & $15$ & $10$ & $1.28358$ & 2 ext. \\
$\mathtt{uuvuuvvuvvuvv}$ & $(13,6)$ & $15$ & $10$ & $1.28358$ & 2 ext. \\
$\mathtt{uuvuvuuvvvv}$ & $(11,5)$ & $14$ & $10$ & $1.28358$ & 2 ext. \\
$\mathtt{uuuvvvuvuvvv}$ & $(12,5)$ & $14$ & $10$ & $1.28358$ & 2 ext. \\
$\mathtt{uuvvuuvvvuvvv}$ & $(13,5)$ & $14$ & $10$ & $1.28358$ & 2 ext. \\
$\mathtt{uuvuvuuvvvvvvv}$ & $(14,5)$ & $26$ & $22$ & $1.28915$ & 5 ext. \\
$\mathtt{uuuuuvuvvvuv}$ & $(12,7)$ & $20$ & $12$ & $1.30443$ & 2 ext. \\
$\mathtt{uuuvuvvvvvuv}$ & $(12,5)$ & $18$ & $12$ & $1.30443$ & 2 ext. \\
$\mathtt{uuuvuvvvvvvvuv}$ & $(14,5)$ & $26$ & $20$ & $1.30499$ & 4 ext. \\
\bottomrule
\end{tabular}
\caption{The $24$ coprime classes of smallest Mahler measure $>1$: word, type, degree of
$P_w$, degree of its non-cyclotomic factor, Mahler measure, and the nature of that factor
(Lehmer's polynomial; a Salem number $=$ reciprocal with one root off the unit circle; a
Pisot number; or ``$k$ ext.'' $=$ $k$ roots outside the circle).}\label{tab:highlights}
\end{table}


\begin{rem}
The cyclotomic part of $P_w$ collects roots of unity; the non-cyclotomic part carries the
arithmetic. Small Mahler measure means a Salem or Pisot factor, and these recur across many
words and types: within $E+o\le20$ Lehmer's value $1.17628$ is realised by several words, and
the next measures attained are $1.21639,\,1.23526,\,1.25105,\ldots$ Specialising $t$ further to
$(X,Y)=(2,3)$ feeds exactly these factors into
the divisibility ceilings of \texttt{cocycle-combine}: a Salem/Lehmer factor is a
\emph{near-cyclotomic} obstruction, with Mahler measure just above $1$. Extending the catalogue
past $E+o=20$, and pinning down which words realise the successive small Salem numbers, is the
natural continuation.
\end{rem}

\subsection*{The map $w\mapsto R_w$}

The core is a true invariant of the cyclic word, with two reflection symmetries. Write $w^{R}$
for the reversed word and $\bar w$ for the complement (swap $u\leftrightarrow v$, which sends
type $(E,o)$ to $(E,E-o)$).

\begin{prop}[Reflections and fibres]\label{prop:cores}
$R_{w^{R}}$ and $R_{\bar w}$ both equal the reciprocal $t^{\deg R_w}R_w(1/t)$. Hence
$R_{\overline{w^{R}}}=R_w$: the involution $w\mapsto\overline{w^{R}}$, of type
$(E,o)\mapsto(E,E-o)$, preserves the core, so the cores occur in fibres closed under it. Moreover
$R_w(0)=R_w(1)=1$.
\end{prop}

\begin{proof}
$R_w(0)=P_w(0)/g_o(0)=1$ (both have constant term $1$) and $R_w(1)=1$ by
Corollary~\ref{cor:Rone}. Reversal exchanges prefixes with suffixes and complement exchanges
$(o,E-o)$; on the reduced cocycle each acts by $t\mapsto t^{-1}$ --- for the complement via the
Fox fundamental formula $(X-1)S_u(w)+(Y-1)S_v(w)=\chi(w)-1$ together with $S_u(\bar w)=S_v(w)$
--- and $g_o$ is reciprocal. Verified on all coprime necklaces with $E\le16$.
\end{proof}

The map is very far from injective: the $4506$ coprime necklaces with $E\le16$ realise only
$1529$ distinct cores ($818$ up to reciprocal). The simple cores have enormous fibres ---
$R_w=1$ for $80$ necklaces, $R_w=\Phi_6=t^2-t+1$ for $50$ --- and the $55$ \emph{cyclotomic}
cores (Mahler $1$) have mean fibre $7.7$ against $2.8$ for the $1474$ non-cyclotomic ones. So the
arithmetic of a word concentrates in a rare non-cyclotomic core shared by only a few necklaces.
Table~\ref{tab:cores} lists the small-degree cores; the full distinct-core catalogue is
\texttt{R\_cores.csv}.

\begin{table}[h]\centering\small\renewcommand{\arraystretch}{1.2}
\begin{tabular}{@{}lcccc@{}}
\toprule
core $R_w(t)$ & $\deg$ & $\mathrm{M}$ & cyclotomic & fibre\\
\midrule
$1$ & $0$ & $1$ & yes & $80$\\
$t^{2} - t + 1$ & $2$ & $1$ & yes & $50$\\
$t^{3} - t^{2} + 1$ & $3$ & $1.3247$ & no & $38$\\
$t^{3} - t + 1$ & $3$ & $1.3247$ & no & $38$\\
$t^{4} - t^{3} + t^{2} - t + 1$ & $4$ & $1$ & yes & $34$\\
$t^{4} - t^{3} + 1$ & $4$ & $1.4013$ & no & $24$\\
$t^{4} - t^{2} + 1$ & $4$ & $1$ & yes & $24$\\
$t^{4} - t + 1$ & $4$ & $1.4013$ & no & $24$\\
\bottomrule
\end{tabular}
\caption{The distinct cores $R_w$ of degree $\le4$, over all $4506$ coprime necklaces with $E\le16$ ($1529$ distinct cores, $55$ of them cyclotomic). Each recurs across a whole fibre of necklaces; cyclotomic cores dominate. Cores come in reciprocal pairs (e.g.\ $t^3-t^2+1$ and $t^3-t+1$), self-reciprocal ones single.}\label{tab:cores}
\end{table}


\paragraph{How $R_w$ splits.} For coprime type the ambient ring $\Lam/(1-\chi(w))=\Z[t,t^{-1}]$
is a domain, so ``splitting'' is the factorisation of $R_w$ into cyclotomic factors times a
non-cyclotomic Salem/Pisot factor --- the latter absent exactly for the $55$ cyclotomic cores.
For \emph{non-coprime} type, $d=\gcd(E,o)>1$, the ring itself splits: by the Chinese remainder
theorem (over $\Q$)
\[
   \Q[t,t^{-1},s]/(s^{d}-1)\;\cong\;\prod_{e\mid d}\Q(\zeta_e)[t,t^{-1}],
\]
one factor per divisor $e$ of $d$, and the reduced cocycle distributes over the factors; the
$e=1$ component recovers the present $R_w$, while the $e>1$ components carry genuinely new
$\zeta_e$-twisted data. The non-coprime stratum is the natural next target.
